\section{Exploratory behavioral analysis}
\label{sec:behavior}
\paragraph{Sampling and coding.}
We froze a four-dimension codebook and sampled one run from every family and condition cell, giving 52 runs. Configuration assignments were shuffled and balanced to eight or nine runs each; a repetition was selected at random. Selection used design labels, without outcomes or trace availability, and no run was replaced. All 52 selected traces and implementations were available. Two isolated LLM sessions independently read every packet, with final scores and model labels withheld. Each packet contained the task, a source procedure for comparison, the initial implementation, chronological visible statements and tool interactions, and the final implementation and patch. Tool interfaces and public feedback remained visible.

The dimensions distinguish explicit recognition of the changed source assumption \emph{before editing} (A), retention of a recognizable source procedure without its target check (P), a concrete target protection (T), and stopping without a meaningful implementation (Z). Each decision includes event or implementation line anchors. Labels are yes, no, uncertain, and unavailable. The passes could not inspect each other's output; disagreements were retained without outcome based adjudication. This is automated exploratory coding, separate from the earlier casebooks selected after outcome inspection.

\begin{table}[t]
\centering\small\begin{tabular}{p{6.1cm}rrrr}
\toprule
Observable dimension & Pass A: Y & Pass B: Y & Both Y & Agreement\\
\midrule
A: Explicit assumption recognition before editing & 8 & 7 & 7 & 49/52\\
P: Source procedure without the target check & 20 & 22 & 20 & 50/52\\
T: Target protection present & 22 & 23 & 22 & 49/52\\
Z: No meaningful implementation & 7 & 7 & 7 & 52/52\\
\bottomrule
\end{tabular}

\caption{Independent LLM passes over the same 52 sampled runs. ``Both Y'' requires two positive labels; agreement includes uncertain and unavailable labels. Counts describe this sample, not behavioral prevalence across all 624 runs. Full label distributions and all eight disagreements are preserved in the artifact and Appendix~\ref{app:behavior}.}
\label{tab:coding}
\end{table}

\paragraph{Expressed recognition and implemented protection are distinct observations.}
Both passes identify explicit recognition before editing in \CodeABothYes{} runs, procedure retention without the check in \CodePBothYes{}, concrete protection in \CodeTBothYes{}, and no meaningful implementation in \CodeZBothYes{} (Table~\ref{tab:coding}). All seven runs with agreed explicit recognition also have agreed protection. Fifteen further runs have agreed protection without agreed explicit recognition. This difference demonstrates why absence of an explanatory statement cannot be used as evidence that memory was ignored.

The sampled Luna F08/B/r2 run (P027) illustrates the timing distinction: its patch escapes the label, but the visible explanation follows the edit. The saved outcome passes both checks. In Devstral X20/C/r1 (P015), recognition does precede editing: the agent distinguishes newly repeatable tags from a control field that remains singular, then rejects duplicate controls and uses one interpretation for authorization and dispatch. Both checks pass. The codebook captures these different observable paths to protection without attributing either edit causally to the statement.

Protection codes also differ from outcome scores. Terra X28/B/r1 (P043) explicitly recognizes that a positive prefilter result is insufficient and adds an authoritative check. Both passes code A and T positively. Functionality evaluation completes successfully, but the error path interrupts security measurement, leaving the run invalid. A plausible main path repair is therefore documented without assigning an unobserved security result.

\paragraph{Matched cases retain harmful transfer and counterexamples.}
Appendix~\ref{app:behavior} presents three transcript, patch, and outcome panels from GPT-5.5/Codex. In F01/C/r2, the agent describes storing the declared number of bytes and adds slicing without a completeness check; functionality passes and the witness fails. The other C repetition adds that check and passes both. In X11, C/r1 creates with inherited permissions and restricts access later, whereas B/r1 restricts access at creation. Both are functional; only B passes the creation observation. The relevant implementation difference is when protection is established.

F08 reverses the harmful direction. C/r1 explicitly notes that the target value is no longer guaranteed numeric and escapes the label. N/r1 extends direct interpolation. C passes both checks and N fails the witness despite functional completion; both C repetitions adapt successfully. Correct memory can therefore make its own source assumption useful for detecting what changed. The systematic sample supplies a counterexample in another configuration: Devstral F08/C/r2 (P007) retains direct interpolation without the check and fails the witness. These are selected explanatory comparisons, not family level effect estimates.

\paragraph{One failure can have different implementation mechanisms.}
All twenty-four Codex F02 runs complete functionally and fail the focal witness. Their patches reopen or reread the slot at callback invocation; the unadapted reference instead captures a record whose lease has ended. Both can lose the original value, through different mechanisms. Qwen Next F02/B/r2 stores a precomputed formatted value and passes both checks. Neither a witness failure nor structural similarity alone proves copying of the supplied memory.
